% Options for packages loaded elsewhere
\PassOptionsToPackage{unicode}{hyperref}
\PassOptionsToPackage{hyphens}{url}
\PassOptionsToPackage{dvipsnames,svgnames,x11names}{xcolor}
\documentclass[
  10pt,
  fontset=fandol]{ctexart}
\usepackage{xcolor}
\usepackage[margin=16mm]{geometry}
\usepackage{amsmath,amssymb}
\setcounter{secnumdepth}{-\maxdimen} % remove section numbering
\usepackage{iftex}
\ifPDFTeX
  \usepackage[T1]{fontenc}
  \usepackage[utf8]{inputenc}
  \usepackage{textcomp} % provide euro and other symbols
\else % if luatex or xetex
  \usepackage{unicode-math} % this also loads fontspec
  \defaultfontfeatures{Scale=MatchLowercase}
  \defaultfontfeatures[\rmfamily]{Ligatures=TeX,Scale=1}
\fi
\usepackage{lmodern}
\ifPDFTeX\else
  % xetex/luatex font selection
\fi
% Use upquote if available, for straight quotes in verbatim environments
\IfFileExists{upquote.sty}{\usepackage{upquote}}{}
\IfFileExists{microtype.sty}{% use microtype if available
  \usepackage[]{microtype}
  \UseMicrotypeSet[protrusion]{basicmath} % disable protrusion for tt fonts
}{}
\makeatletter
\@ifundefined{KOMAClassName}{% if non-KOMA class
  \IfFileExists{parskip.sty}{%
    \usepackage{parskip}
  }{% else
    \setlength{\parindent}{0pt}
    \setlength{\parskip}{6pt plus 2pt minus 1pt}}
}{% if KOMA class
  \KOMAoptions{parskip=half}}
\makeatother
\setlength{\emergencystretch}{3em} % prevent overfull lines
\providecommand{\tightlist}{%
  \setlength{\itemsep}{0pt}\setlength{\parskip}{0pt}}
\usepackage{amsmath,amssymb,mathtools,needspace}
\xeCJKDeclareCharClass{CJK}{"2032,"0400->"04FF,"2160->"216B,"2460->"2473,"25A0,"25C7,"25CB,"25CF}
\IfFontExistsTF{Arial Unicode MS}{\xeCJKsetup{AutoFallBack=true}\setCJKfallbackfamilyfont{\CJKrmdefault}{Arial Unicode MS}}{}
\setcounter{secnumdepth}{0}
\setlength{\emergencystretch}{2em}
\allowdisplaybreaks
\usepackage{bookmark}
\IfFileExists{xurl.sty}{\usepackage{xurl}}{} % add URL line breaks if available
\urlstyle{same}
\hypersetup{
  pdftitle={算法设计关系到科学计算的成败},
  colorlinks=true,
  linkcolor={Maroon},
  filecolor={Maroon},
  citecolor={Blue},
  urlcolor={Blue},
  pdfcreator={LaTeX via pandoc}}

\title{算法设计关系到科学计算的成败}
\author{}
\date{}

\begin{document}
\maketitle

\subsubsection{13.1.1 算法设计关系到科学计算的成败}\label{ux7b97ux6cd5ux8bbeux8ba1ux5173ux7cfbux5230ux79d1ux5b66ux8ba1ux7b97ux7684ux6210ux8d25}

计算机是一种功能很强的计算工具。现代超级计算机的运算速度已高达每秒亿亿次。计算机运算速度如此之快，是否意味着计算机上的算法可以随意选择呢？

举个简单的例子。

众所周知，Cramer 法则原则上可用来求解线性方程组。用这种方法求解一个 \(n\) 阶方程组，要计算 \(n+1\) 个 \(n\) 阶行列式的值，总共要做 \((n+1)n!(n-1)\) 次乘除操作。当 \(n\) 充分大时，这个计算量是惊人的。譬如一个不算太大的 20 阶线性方程组，大约要做 \(10^{21}\) 次乘除操作，这项计算即使用每秒 30 万亿次的超级计算机来承担，也得要连续工作

\[
\frac{10^{21}}{3\times10^{11}\times60\times60\times24\times365}\approx1\text{（年）}
\]

才能完成。当然这是完全没有实际意义的。

其实，求解线性方程组有许多实用解法。譬如，运用人们熟悉的消元技术，一个 20 阶的线性方程组即使用普通的计算器也能很快地解出来。这个简单的例子说明，\textbf{能否合理地选择算法是科学计算成败的关键。}

随着计算机的广泛应用与日益普及，算法设计的重要性正越来越为人们所认识。《计算机大百科全书》在其``算法学''词条中指出：``凡与计算机打交道，无不研究各种类型的算法。''\textbf{``算法学是计算机科学最重要的内容，有的计算机学者甚至称，计算机科学就是算法的科学。''}

\begin{quote}
校注（agent 补充，CH13-E001）：原页正文印``每秒 30 万亿次''，上式分母却印 \(3\times10^{11}\)，且计算结果印为``\(\approx1\)（年）''。此处按原式保留。\(30\) 万亿 \(=3\times10^{13}\)，用正文速度计算约为 \(1.057\) 年；用原式分母计算约为 \(105.7\) 年。疑为原书分母指数排印错误，非 OCR 错误。已对照放大原图核对指数为 \(11\)。
\end{quote}

\subsection{本节覆盖原页的校注记录}\label{ux672cux8282ux8986ux76d6ux539fux9875ux7684ux6821ux6ce8ux8bb0ux5f55}

下列记录按原页关联，含同页相邻小节的校注；计算时同时读取上文校注和完整原页。

\begin{itemize}
\tightlist
\item
  \texttt{CH13-E001}：原书 PDF 页 307
\end{itemize}

\end{document}
